\section{Transformer Attention：先把比较对象写清楚}

课程没有直接把 SDM 贴到 Transformer 上，而是完整复习了一次 attention。这样做是为了锁定比较对象：query 与 key 的相似度怎样形成权重，softmax 做了什么，value 为什么不同于 key，以及最终更新向量如何写成矩阵式。

\subsection{Q、K、V 的操作语义}

本节先把后续几何比较所依赖的 Q、K、V 操作语义固定下来。章节页之后，语言示例把当前待预测位置视为 query，把上下文 token 的 key 用于匹配，把 value 视为真正被搬运到输出的内容。Q/K/V 页把流程拆成生成投影、query--key 比较、归一化和 value 加权；下一页展示归一化后的权重怎样乘上各 value 再相加。

\lecturefigure{slide-13-transformer-attention-overview.jpg}{Transformer attention 章节过渡}{官方录像恢复幻灯片 V013；00:07:20}

\lecturefigure{slide-14-attention-example.jpg}{语言模型中的 attention 示例}{官方录像恢复幻灯片 V014；00:08:02}

\lecturefigure{slide-15-qkv-steps.jpg}{Keys、queries 与 values 的四步计算}{官方录像恢复幻灯片 V015；00:08:40}

\lecturefigure{slide-16-weighted-value-sum.jpg}{Softmax 权重如何汇总 value vectors}{官方录像恢复幻灯片 V016；00:09:00}

\begin{importantbox}{Key 与 value 不能混为一谈}
key 回答“这条记忆与 query 有多相关”，value 回答“若相关，应取回什么”。把二者分开后，模型可以按语法或实体身份匹配，却返回对下一 token 有用的不同特征；这也对应 SDM 中 address 与 stored content 可分离的 hetero-association。
\end{importantbox}

\subsection{虚构例子与完整更新式}

Fictional example 页用 ``cat''、``sat'' 一类 token 说明高相关 key 可以返回包含语义与句法信息的 value。Full query update 页随后把投影与更新写在一起。设输入 token 矩阵为 $X$，查询位置向量为 $x_q$，投影矩阵为 $W_Q,W_K,W_V$，则
$$
q=W_Qx_q,\quad K=XW_K,\quad V=XW_V,
$$
$$
\operatorname{Attn}(q,K,V)
=V^{\top}\operatorname{softmax}\!\left(\beta Kq\right).
$$
其中 $q$ 是 query，$K$ 的每一行是 key，$V$ 的每一行是 value，$\beta$ 是温度/尺度系数。若采用 scaled dot-product，常见 $\beta=1/\sqrt{d_k}$。

\lecturefigure{slide-17-fictional-example.jpg}{虚构例子：匹配 key 后返回不同 value 特征}{官方录像恢复幻灯片 V017；00:10:14}

\lecturefigure{slide-18-full-query-update.jpg}{完整 query update 的矩阵表达}{官方录像恢复幻灯片 V018；00:10:30}

\subsection{Dot product、softmax 与 value summation}

三页分别放大 attention 的关键操作。Dot-product 页说明相似度来自 $k_i^{\top}q$；softmax 页说明指数化会扩大高分差异并归一化；summation 页说明输出是 value 的凸组合。Softmax 的关键不是“产生概率”这个标签，而是权重关于相似度呈指数变化，这正是下一章要从高维球交集中恢复的形状。

\lecturefigure{slide-19-query-key-dot-product.jpg}{Query--key dot product 形成 similarity logits}{官方录像恢复幻灯片 V019；00:10:50}

\lecturefigure{slide-20-softmax.jpg}{Softmax 放大高相似项并完成归一化}{官方录像恢复幻灯片 V020；00:11:44}

\lecturefigure{slide-21-summation-over-values.jpg}{按注意力权重求和 value vectors}{官方录像恢复幻灯片 V021；00:12:02}

\lecturefigure{slide-22-full-attention-equation.jpg}{Attention 全流程与最终更新式}{官方录像恢复幻灯片 V022；00:12:18}

对第 $i$ 个 key，softmax 权重为
$$
\alpha_i(q)=\frac{\exp(\beta k_i^{\top}q)}
{\sum_j\exp(\beta k_j^{\top}q)},
\qquad
y(q)=\sum_i\alpha_i(q)v_i.
$$
这里 $\alpha_i$ 是 normalized retrieval weight，$v_i$ 是第 $i$ 条记忆内容。与 SDM 比较时，问题转化为：sphere-intersection count 是否可写成 $\exp(\text{similarity})$，以及归一化后是否得到同样的 $\alpha_i$。

\begin{lstlisting}[caption={Scaled dot-product attention（单 query 示意）}]
def attention(query, keys, values, beta):
    logits = beta * (keys @ query)
    weights = softmax(logits)
    return weights @ values
\end{lstlisting}

\teachervoice{讲者特别提醒“what you pay attention to”和“what you get out”可以不同：query 与 keys 决定相关性，values 决定取回内容。这句话是后面把 cerebellar activation pathway 与 write pathway 分开的关键，也连接了第二讲中的 content/structure factorization。}

\subsection{本章小结}

Attention 是内容寻址记忆：query 与 keys 产生相似度，softmax 将相似度变为指数归一化权重，values 被加权取回。只要 SDM 的 sphere intersection 随距离指数衰减，并能把距离改写成点积，两种读操作就拥有相同的计算骨架。

